%% ============================================================
%% Supplementary Material
%% Multi-paradigm benchmarking for integrated soil quality
%% assessment under conservation management in a tropical Ultisol
%% ============================================================
\documentclass[12pt]{article}
\usepackage[a4paper,margin=2.5cm]{geometry}
\usepackage{booktabs}
\usepackage{multirow}
\usepackage{longtable}
\usepackage{array}
\usepackage{amsmath}
\usepackage{caption}
\usepackage[hidelinks]{hyperref}

\captionsetup{labelfont=bf, font=small}
\renewcommand{\thetable}{S\arabic{table}}

\begin{document}

\begin{center}
{\Large\bfseries Supplementary Material}\\[6pt]
{\large Multi-paradigm benchmarking for integrated soil quality assessment\\
under conservation management in a tropical Ultisol}\\[12pt]
\end{center}

\noindent This document contains four supplementary tables referenced in the main manuscript.

%% ============================================================
%%  TABLE S1 — PCA eigenvalues and component weights
%% ============================================================
\clearpage
\begin{table}[ht]
\centering
\caption{PCA eigenvalues, variance explained, and component weights for the ANOVA/MDS paradigm (34 indicators). Only the eight retained components (eigenvalue $> 1$) are shown; together they explain 83.56\% of total variance.}\label{tab:s1_eigen}
\begin{tabular}{lcccc}
\toprule
Component & Eigenvalue & Var.\ expl.\ (\%) & Cumul.\ (\%) & Weight$^a$ \\
\midrule
PC1  & 8.230 & 24.21 & 24.21 & 0.310 \\
PC2  & 5.695 & 16.75 & 40.96 & 0.215 \\
PC3  & 3.284 & 9.66  & 50.61 & 0.124 \\
PC4  & 2.934 & 8.63  & 59.24 & 0.111 \\
PC5  & 2.566 & 7.55  & 66.79 & 0.097 \\
PC6  & 2.435 & 7.16  & 73.95 & 0.092 \\
PC7  & 1.794 & 5.28  & 79.23 & 0.068 \\
PC8  & 1.471 & 4.33  & 83.56 & 0.056 \\
\bottomrule
\end{tabular}
\smallskip

\noindent{\footnotesize $^a$Weight = eigenvalue$_i$ / $\sum_{i=1}^{8}$eigenvalue$_i$; used in the additive SQI$_{\text{ANOVA}}$ computation.}
\end{table}

%% ============================================================
%%  TABLE S2 — Complete fuzzy rule base
%% ============================================================
\clearpage
\begin{table}[ht]
\centering
\caption{Complete expert rule base of the hierarchical Mamdani fuzzy inference system (IPACSA). Membership functions (MFs): L\,=\,Low, M\,=\,Medium, H\,=\,High; 0\,=\,don't care. Connection: AND (minimum t-norm). Defuzzification: centroid.}\label{tab:s2_fuzzy}
\smallskip
\begin{tabular}{p{4.5cm}lp{5cm}cc}
\toprule
Sub-index (inputs) & Rule & Antecedent pattern & Output & Weight \\
\midrule
\multicolumn{5}{l}{\textit{COS -- Soil Conservation (7 inputs: DMG, DMP, RMP, Ds, Na, ICV, COS$_{\text{stock}}$)}} \\
 & E1 & M, M, L, L, L, H, H & H & 1.0 \\
 & E2 & M, M, M, M, M, M, M & M & 1.0 \\
 & E3 & L, L, H, H, H, L, L & L & 1.0 \\
 & F1--F21 & Main-effect fallback$^a$ & L/M/H & 0.1 \\
\midrule
\multicolumn{5}{l}{\textit{ARQ -- Plant Architecture (6 inputs: ALT, DESP, CESP, NESP, NESPC, PESP)}} \\
 & E1 & H, H, H, H, H, H & H & 1.0 \\
 & E2 & M, M, M, M, M, M & M & 1.0 \\
 & E3 & L, L, L, L, L, L & L & 1.0 \\
 & F1--F18 & Main-effect fallback$^a$ & L/M/H & 0.1 \\
\midrule
\multicolumn{5}{l}{\textit{STAND -- Plant Stand (1 input: STAND)}} \\
 & E1 & H & H & 1.0 \\
 & E2 & M & M & 1.0 \\
 & E3 & L & L & 1.0 \\
\midrule
\multicolumn{5}{l}{\textit{IPACSA -- Final aggregation (4 inputs: COS, ARQ, STAND, PROD)}} \\
 & E1 & M, H, M, H & H & 1.20 \\
 & E2 & M, M, M, M & M & 1.00 \\
 & E3 & L, L, L, L & L & 0.80 \\
 & E4 & H, M, M, M & H & 1.10 \\
 & E5 & M, H, M, H & H & 1.10 \\
 & E6 & M, M, H, M & H & 1.10 \\
 & E7 & M, H, M, M & H & 1.15 \\
 & F1--F12 & Main-effect fallback$^a$ & L/M/H & 0.1 \\
\bottomrule
\end{tabular}
\smallskip

\noindent{\footnotesize $^a$Main-effect fallback rules: for each input variable $k$, three rules are generated where only the $k$-th input is set to L, M, or H (all others = 0, i.e., don't care), mapping directly to the corresponding output level. These rules (weight\,=\,0.1) prevent NaN outputs in regions of the input space not covered by expert rules.}
\end{table}

%% ============================================================
%%  TABLE S3 — PLS-SEM diagnostics
%% ============================================================
\clearpage

\begin{table}[ht]
\centering
\caption{PLS-SEM structural model: path coefficients, bootstrap confidence intervals, and model fit (five constructs, 25 manifest variables after VIF screening).}\label{tab:s3_paths}
\begin{tabular}{lcccccc}
\toprule
Path & $\gamma$ & Boot.\ Mean & Boot.\ SD & $t$ & 95\% CI & $p$ \\
\midrule
PHYSICAL $\to$ YIELD  & $-$0.122 & $-$0.013 & 0.160 & 0.76  & [$-$0.28,\;0.25]  & 0.992 \\
CHEMICAL $\to$ YIELD  &    0.097 &    0.104 & 0.102 & 0.95  & [$-$0.14,\;0.28]  & 0.264 \\
MICRO $\to$ YIELD     &    0.175 &    0.119 & 0.127 & 1.39  & [$-$0.22,\;0.28]  & 0.284 \\
AGRO $\to$ YIELD      &    0.622 &    0.608 & 0.051 & 12.30 & [\phantom{$-$}0.50,\;0.70] & $<$0.001 \\
\midrule
\multicolumn{7}{l}{$R^2 = 0.690$; Adj.\ $R^2 = 0.678$; Bootstrap: 500 resamples, seed = 42.} \\
\bottomrule
\end{tabular}
\end{table}

\begin{table}[ht]
\centering
\caption{PLS-SEM outer weights (Mode~B, formative) and bootstrap significance for 25 retained manifest variables.}\label{tab:s3_weights}
\begin{tabular}{llrrrrl}
\toprule
Construct & Indicator & Weight & Boot.\ SD & $t$ & 95\% CI & $p$ \\
\midrule
\multirow{9}{*}{PHYSICAL}
 & DMG   &  0.068 & 0.192 & 0.35 & [$-$0.29,\;0.41] & 0.792 \\
 & COS   &  0.086 & 0.217 & 0.40 & [$-$0.39,\;0.47] & 0.784 \\
 & Na    &  0.153 & 0.216 & 0.71 & [$-$0.36,\;0.43] & 0.864 \\
 & ICV   &  0.031 & 0.196 & 0.16 & [$-$0.34,\;0.44] & 0.888 \\
 & MI    & $-$0.301 & 0.315 & 0.95 & [$-$0.50,\;0.60] & 0.800 \\
 & MA    &  0.684 & 0.691 & 0.99 & [$-$1.09,\;1.17] & 0.924 \\
 & VIB   &  0.178 & 0.268 & 0.66 & [$-$0.48,\;0.55] & 0.900 \\
 & RP    &  0.314 & 0.374 & 0.84 & [$-$0.66,\;0.68] & 0.980 \\
 & AD/PT &  1.200 & 1.108 & 1.08 & [$-$1.41,\;1.47] & 0.924 \\
\midrule
\multirow{6}{*}{CHEMICAL}
 & pH    & $-$0.365 & 0.206 & 1.77 & [$-$0.61,\;0.42] & 0.092 \\
 & P     & $-$0.508 & 0.263 & 1.93 & [$-$0.75,\;0.52] & 0.088 \\
 & K     & $-$0.044 & 0.242 & 0.18 & [$-$0.51,\;0.44] & 0.936 \\
 & Mg    &  0.096 & 0.196 & 0.49 & [$-$0.32,\;0.44] & 0.640 \\
 & CEC   & $-$0.197 & 0.210 & 0.94 & [$-$0.51,\;0.46] & 0.316 \\
 & EstN  &  0.618 & 0.330 & 1.87 & [$-$0.51,\;1.06] & 0.088 \\
\midrule
\multirow{4}{*}{MICRO}
 & qMIC  & $-$0.580 & 0.445 & 1.30 & [$-$0.98,\;0.80] & 0.276 \\
 & qCO$_2$ &  0.357 & 0.295 & 1.21 & [$-$0.43,\;0.78] & 0.276 \\
 & CCO$_2$ & $-$0.246 & 0.363 & 0.68 & [$-$0.72,\;0.75] & 0.656 \\
 & NMIC  &  0.927 & 0.596 & 1.56 & [$-$0.90,\;1.05] & 0.276 \\
\midrule
\multirow{4}{*}{AGRO}
 & ALT   &  0.601 & 0.089 & 6.74 & [\phantom{$-$}0.40,\;0.76] & $<$0.001 \\
 & DESP  &  0.124 & 0.087 & 1.43 & [$-$0.06,\;0.30] & 0.172 \\
 & CESP  &  0.337 & 0.081 & 4.16 & [\phantom{$-$}0.17,\;0.48] & $<$0.001 \\
 & STAND &  0.529 & 0.108 & 4.90 & [\phantom{$-$}0.32,\;0.72] & 0.004 \\
\midrule
\multirow{2}{*}{YIELD}
 & NESP  &  0.534 & 0.051 & 10.46 & [\phantom{$-$}0.44,\;0.62] & $<$0.001 \\
 & PESP  &  0.643 & 0.052 & 12.34 & [\phantom{$-$}0.56,\;0.76] & $<$0.001 \\
\bottomrule
\end{tabular}
\smallskip

\noindent{\footnotesize VIF screening (threshold = 5) excluded Ca (VIF\,=\,26.5), MOS (VIF\,=\,9.8), MBC (VIF\,=\,23.9), and GY (VIF\,=\,20.1). YIELD construct estimated with Mode~A (reflective); all others with Mode~B (formative). Bootstrap: 500 resamples.}
\end{table}

%% ============================================================
%%  TABLE S4 — RF variable importance
%% ============================================================
\clearpage
\begin{table}[ht]
\centering
\caption{Random Forest variable importance (34-indicator model, 500 trees, $m_{\text{try}} = 11$). Variables ranked by permutation importance (\%IncMSE).}\label{tab:s4_rf}
\begin{tabular}{rlccc}
\toprule
Rank & Indicator & Domain & \%IncMSE & IncNodePurity \\
\midrule
1  & NESP  & Yield        & 20.22 & 30.46 \\
2  & STAND & Agronomic    & 19.49 & 23.51 \\
3  & NESPC & Yield        & 13.19 & 8.44 \\
4  & PROD  & Yield        & 13.15 & 8.02 \\
5  & PESP  & Yield        & 12.51 & 8.26 \\
6  & DESP  & Agronomic    & 12.33 & 9.38 \\
7  & ICV   & Physical     & 11.19 & 5.95 \\
8  & CESP  & Agronomic    & 11.13 & 5.27 \\
9  & P     & Chemical     & 8.04 & 3.01 \\
10 & Ca    & Chemical     & 7.63 & 1.32 \\
11 & SOM   & Chemical     & 7.35 & 1.68 \\
12 & CEC   & Chemical     & 6.68 & 1.48 \\
13 & K     & Chemical     & 6.64 & 1.81 \\
14 & MI    & Physical     & 6.47 & 1.16 \\
15 & pH    & Chemical     & 6.38 & 1.35 \\
16 & EstN  & Chemical     & 6.11 & 1.36 \\
17 & Al    & Chemical     & 6.09 & 1.23 \\
18 & qMIC  & Microbiol.   & 5.89 & 1.11 \\
19 & MBC   & Microbiol.   & 5.64 & 0.97 \\
20 & AD/PT & Physical     & 5.51 & 1.20 \\
21 & VIB   & Physical     & 5.39 & 1.16 \\
22 & BR    & Microbiol.   & 5.23 & 1.12 \\
23 & MBN   & Microbiol.   & 5.09 & 1.15 \\
24 & RP    & Physical     & 4.96 & 0.86 \\
25 & Mg    & Chemical     & 4.68 & 0.60 \\
26 & RMP   & Physical     & 4.53 & 1.29 \\
27 & qCO$_2$  & Microbiol. & 4.38 & 0.80 \\
28 & MA    & Physical     & 4.31 & 1.00 \\
29 & COS   & Physical     & 4.18 & 1.78 \\
30 & GMD   & Physical     & 3.87 & 1.10 \\
31 & ALT   & Agronomic    & 3.40 & 1.39 \\
32 & WMD   & Physical     & 3.30 & 1.56 \\
33 & Na    & Physical     & 2.97 & 1.10 \\
34 & BD    & Physical     & 2.95 & 1.05 \\
\bottomrule
\end{tabular}
\smallskip

\noindent{\footnotesize \%IncMSE, percentage increase in mean squared error upon permutation of the variable; IncNodePurity, total decrease in node impurity (residual sum of squares). Domain abbreviations: Microbiol.\,=\,Microbiological.}
\end{table}

\end{document}
